\begin{proof}
For the upper bound take $Q=\mathcal L_\ell$, initialize at
$v_0=0^{\lceil\ell/2\rceil}$, update by
\[
 U_a(v)=\operatorname{trim}_\ell(av),\qquad
 D(v)=\tau_v(1/2),
\]
where trimming retains the shortest prefix of cost at least $\ell$.
Such a prefix always exists in $av$. Initially $x_0\in I_{v_0}$.
If $x_t\in I_v$ and $av=v'u$, invariance of $[0,1]$ implies
$x_{t+1}\in\tau_{v'}(\tau_u([0,1]))\subseteq I_{v'}$.
Induction proves this inclusion for all histories. The midpoint readout
therefore has error at most $|I_{v'}|/2\le\theta/2$, including startup.
Every leaf is reachable by supplying its letters in chronological order,
since $\operatorname{trim}_\ell(vv_0)=v$.

For the lower bound one must use reachable histories, rather than assume
that every point of $K$ occurs at finite time. Set $T=\ell+3$.
For each leaf $v$ first supply $T-|v|$ copies of letter 1 and then the
letters of $v$ in reverse order. The final true state is
\[
 y_v=\tau_v\bigl(\tau_1^{T-|v|}(0)\bigr),\qquad
 |b_v-y_v|=r_v2^{-(T-|v|)}\le\theta/8,
\]
because $|v|\le\ell$. Hence two different $y_v$ are separated by at
least $3\theta/2-2\theta/8=5\theta/4>\theta$.
All these histories have the same length. If two ended in the same
stored state, their common readout would be within $\theta/2$ of both,
contradicting this separation. Thus at least $F_{\ell+2}$ states are
necessary. Finally
$F_{130}=659034621587630041982498215$ lies strictly between $2^{89}$
and $2^{90}$, proving the numerical bit optimum.
\end{proof}
